\section{Ergebnisse und Evaluation}

Dieses Kapitel stellt die Ergebnisse der prototypischen Implementierung vor und bewertet diese anhand der definierten Anforderungen.

\subsection{Funktionstest}

Zur Validierung wurde ein manueller Funktionstest durchgeführt, der den vollständigen Prozess von der Anlage einer Bestellanforderung bis zum Abschluss des Genehmigungsworkflows abdeckt.

Der \textbf{Genehmigungspfad} umfasst: Anlegen einer Bestellanforderung im Status \texttt{Draft}, Erfassen von Kopfdaten und Positionen, automatische Betragsberechnung beim Speichern, Freigabe über die \textit{Freigeben}-Aktion, Start des Workflows in \gls{spa}, Genehmigung in der My Inbox und Statusaktualisierung auf \texttt{Approved} über den Callback.

Der \textbf{Ablehnungspfad} verläuft analog, wobei der Status auf \texttt{Rejected} gesetzt wird und eine erneute Bearbeitung und Einreichung möglich ist.

Tabelle~\ref{tab:testfaelle} fasst die Testfälle zusammen.

\begin{table}[H]
\centering
\caption{Durchgeführte Testfälle}
\label{tab:testfaelle}
\begin{tabularx}{\textwidth}{c|X|c}
\hline
\textbf{ID} & \textbf{Testfall} & \textbf{Ergebnis} \\
\hline
T1 & Bestellanforderung anlegen und speichern & Bestanden \\
T2 & Automatische Betragsberechnung & Bestanden \\
T3 & Freigeben ohne Positionen (Validierung) & Bestanden \\
T4 & Freigeben mit Positionen, Workflow-Start & Bestanden \\
T5 & Genehmigung in SPA My Inbox & Bestanden \\
T6 & Ablehnung in SPA My Inbox & Bestanden \\
T7 & Draft-Unterstützung & Bestanden \\
T8 & Dokumentenanhänge & Bestanden \\
T9 & Statushistorie & Bestanden \\
T10 & Wiederbearbeitung nach Entscheidung & Bestanden \\
\hline
\end{tabularx}
\end{table}

\subsection{Anforderungsabgleich}

Alle elf funktionalen Anforderungen aus Kapitel~3 wurden vollständig umgesetzt: \gls{crud}-Operationen (FA1), Positionserfassung mit Betragsberechnung (FA2, FA3), Statusverwaltung (FA4), Freigabeprozess mit Validierung (FA5), Workflow-Integration (FA6, FA7), Draft-Unterstützung (FA8), Statushistorie (FA9), Dokumentenanhänge (FA10) und Wiederbearbeitbarkeit (FA11).

\subsection{Bewertung}

\paragraph{Stärken}
Die Kombination von \gls{cap} und Fiori Elements ermöglicht eine hohe Entwicklungsgeschwindigkeit durch deklarative Modellierung. Die Benutzeroberfläche wird vollständig annotationsgesteuert generiert, ohne manuellen Frontend-Code. Die durchgängige Integration innerhalb der \gls{btp} und die klare Schichtentrennung fördern Wartbarkeit und Erweiterbarkeit.

\paragraph{Herausforderungen}
Technische Herausforderungen umfassten die Feldnamen-Konvertierung zwischen \gls{cap} und \gls{spa}, die regionenübergreifende Destination-Konfiguration sowie die Betragsberechnung im Zusammenspiel mit dem Draft-Mechanismus.

\paragraph{Limitierungen}
Der Prototyp ist als Einzelmandantenlösung konzipiert, verfügt über keine mehrstufigen Genehmigungsregeln und keine automatisierten Tests. Für einen produktiven Einsatz wären diese Erweiterungen erforderlich.
